% GENERATED FILE. Edit canonical lessons, then run npm run generate:books.
% canonical-source-hash: fnv1a64:367f25418da86d3c
% canonical-lessons: GE-C299-ehe, GE-C299-ehemann, GE-C299-ehefrau, GE-C299-trauung, GE-C299-standesamt

\chapter{Die Ehe}
\label{ch:ge-die-ehe}
\hlchaptermodality{299}

\begin{chapteropening}
\textbf{By the end of this chapter:} \emph{I can talk about marriage: a marriage, a husband, a wife, the wedding ceremony and the registry office.}

Complete the last lesson of chapter 299: I can talk about marriage: a marriage, a husband, a wife, the wedding ceremony and the registry office.
\end{chapteropening}

\section[die Ehe]{die Ehe — the marriage}

\label{lesson:GE-C299-ehe}



\begin{quote}
Before the new one: say the German for the bridegroom, then the German for the bride and groom.
\end{quote}

\subsection*{You'll want to know: die Ehe}
\textbf{die Ehe} — \textquotedblleft{}the marriage\textquotedblright{}.

\begin{quote}
Ihre Ehe hielt fünfzig Jahre.
\end{quote}

\begin{quote}
\emph{Their marriage lasted fifty years.}
\end{quote}

Plural: \emph{die Ehen}. The adjective is \emph{verheiratet}.

\begin{cousinweb}
\emph{Ehe} goes back to Middle High German \emph{ē} (EtymDB), usually traced to Old High German \emph{ēwa} \textquotedblleft{}law, eternity\textquotedblright{}: marriage as a lasting bond.
\end{cousinweb}

\begin{grammarlens}[title={What you've built}]
One more word for telling, explaining and describing at B1.
\end{grammarlens}

\subsection*{Your turn}
\begin{itemize}
\raggedright
  \item \emph{Say it:} \emph{die Ehe}
  \item \emph{Say it:} \emph{die Ehe}, once more
  \item \emph{Say it:} a sentence of your own with \emph{die Ehe}
  \item \emph{Recall:} say the German for the bridegroom, then the German for the bride and groom, then say \emph{die Ehe} again
\end{itemize}

\begin{tcolorbox}[breakable,colback=teal!4,colframe=teal!35!black,title={Before you move on}]
What does die Ehe mean? (\textquotedblleft{}The marriage\textquotedblright{}.) Say it once more.
\end{tcolorbox}
\section[der Ehemann]{der Ehemann — the husband}

\label{lesson:GE-C299-ehemann}



\begin{quote}
Before the new one: say the German for the bride and groom, then the German for the marriage.
\end{quote}

\subsection*{You'll want to know: der Ehemann}
\textbf{der Ehemann} — \textquotedblleft{}the husband\textquotedblright{}.

\begin{quote}
Ihr Ehemann arbeitet als Koch.
\end{quote}

\begin{quote}
\emph{Her husband works as a cook.}
\end{quote}

Plural: \emph{die Ehemänner}. In speech, \emph{mein Mann} is more common.

\begin{cousinweb}
\emph{Ehemann} joins \emph{Ehe} \textquotedblleft{}marriage\textquotedblright{}, Middle High German \emph{ē} (EtymDB), and the inherited \emph{Mann} \textquotedblleft{}man\textquotedblright{}.
\end{cousinweb}

\begin{grammarlens}[title={What you've built}]
One more word for telling, explaining and describing at B1.
\end{grammarlens}

\subsection*{Your turn}
\begin{itemize}
\raggedright
  \item \emph{Say it:} \emph{der Ehemann}
  \item \emph{Say it:} \emph{der Ehemann}, once more
  \item \emph{Say it:} a sentence of your own with \emph{der Ehemann}
  \item \emph{Recall:} say the German for the bride and groom, then the German for the marriage, then say \emph{der Ehemann} again
\end{itemize}

\begin{tcolorbox}[breakable,colback=teal!4,colframe=teal!35!black,title={Before you move on}]
What does der Ehemann mean? (\textquotedblleft{}The husband\textquotedblright{}.) Say it once more.
\end{tcolorbox}
\section[die Ehefrau]{die Ehefrau — the wife}

\label{lesson:GE-C299-ehefrau}



\begin{quote}
Before the new one: say the German for the marriage, then the German for the husband.
\end{quote}

\subsection*{You'll want to know: die Ehefrau}
\textbf{die Ehefrau} — \textquotedblleft{}the wife\textquotedblright{}.

\begin{quote}
Seine Ehefrau spricht drei Sprachen.
\end{quote}

\begin{quote}
\emph{His wife speaks three languages.}
\end{quote}

Plural: \emph{die Ehefrauen}. In speech, \emph{meine Frau} is more common.

\begin{cousinweb}
\emph{Ehefrau} joins \emph{Ehe} \textquotedblleft{}marriage\textquotedblright{} and \emph{Frau} \textquotedblleft{}woman, wife\textquotedblright{}, from Middle High German \emph{vrouwe} (EtymDB).
\end{cousinweb}

\begin{grammarlens}[title={What you've built}]
One more word for telling, explaining and describing at B1.
\end{grammarlens}

\subsection*{Your turn}
\begin{itemize}
\raggedright
  \item \emph{Say it:} \emph{die Ehefrau}
  \item \emph{Say it:} \emph{die Ehefrau}, once more
  \item \emph{Say it:} a sentence of your own with \emph{die Ehefrau}
  \item \emph{Recall:} say the German for the marriage, then the German for the husband, then say \emph{die Ehefrau} again
\end{itemize}

\begin{tcolorbox}[breakable,colback=teal!4,colframe=teal!35!black,title={Before you move on}]
What does die Ehefrau mean? (\textquotedblleft{}The wife\textquotedblright{}.) Say it once more.
\end{tcolorbox}
\section[die Trauung]{die Trauung — the wedding ceremony}

\label{lesson:GE-C299-trauung}



\begin{quote}
Before the new one: say the German for the husband, then the German for the wife.
\end{quote}

\subsection*{You'll want to know: die Trauung}
\textbf{die Trauung} — \textquotedblleft{}the wedding ceremony\textquotedblright{}.

\begin{quote}
Die Trauung war um elf Uhr.
\end{quote}

\begin{quote}
\emph{The wedding ceremony was at eleven.}
\end{quote}

Plural: \emph{die Trauungen}. In a church or at the town hall.

\begin{cousinweb}
\emph{Trauung} is built on \emph{trauen} \textquotedblleft{}to marry, to trust\textquotedblright{}, from Old High German \emph{trūwēn} (EtymDB): the couple trust themselves to each other.
\end{cousinweb}

\begin{grammarlens}[title={What you've built}]
One more word for telling, explaining and describing at B1.
\end{grammarlens}

\subsection*{Your turn}
\begin{itemize}
\raggedright
  \item \emph{Say it:} \emph{die Trauung}
  \item \emph{Say it:} \emph{die Trauung}, once more
  \item \emph{Say it:} a sentence of your own with \emph{die Trauung}
  \item \emph{Recall:} say the German for the husband, then the German for the wife, then say \emph{die Trauung} again
\end{itemize}

\begin{tcolorbox}[breakable,colback=teal!4,colframe=teal!35!black,title={Before you move on}]
What does die Trauung mean? (\textquotedblleft{}The wedding ceremony\textquotedblright{}.) Say it once more.
\end{tcolorbox}
\section[das Standesamt]{das Standesamt — the registry office}

\label{lesson:GE-C299-standesamt}



\begin{quote}
Before the new one: say the German for the wife, then the German for the wedding ceremony.
\end{quote}

\subsection*{You'll want to know: das Standesamt}
\textbf{das Standesamt} — \textquotedblleft{}the registry office\textquotedblright{}.

\begin{quote}
Sie heirateten zuerst auf dem Standesamt.
\end{quote}

\begin{quote}
\emph{They got married at the registry office first.}
\end{quote}

Plural: \emph{die Standesämter}.

\begin{cousinweb}
\emph{Standesamt} joins \emph{Stand} \textquotedblleft{}standing, status\textquotedblright{} and \emph{Amt} \textquotedblleft{}office\textquotedblright{} with the linking \emph{-es-} (EtymDB): the office for civil status.
\end{cousinweb}

\begin{grammarlens}[title={What you've built}]
One more word for telling, explaining and describing at B1.
\end{grammarlens}

\subsection*{Your turn}
\begin{itemize}
\raggedright
  \item \emph{Say it:} \emph{das Standesamt}
  \item \emph{Say it:} \emph{das Standesamt}, once more
  \item \emph{Say it:} a sentence of your own with \emph{das Standesamt}
  \item \emph{Recall:} say the German for the wife, then the German for the wedding ceremony, then say \emph{das Standesamt} again
\end{itemize}

\begin{tcolorbox}[breakable,colback=teal!4,colframe=teal!35!black,title={Before you move on}]
What does das Standesamt mean? (\textquotedblleft{}The registry office\textquotedblright{}.) Say it once more.
\end{tcolorbox}
